% Source: http://tex.stackexchange.com/a/150903/23931
\documentclass{article}
\usepackage[letterpaper,margin=1in]{geometry}
\usepackage{xcolor}
\usepackage{fancyhdr}
\usepackage{tgschola} % or any other font package you like

\pagestyle{fancy}
\fancyhf{}
\fancyhead[C]{%
  \footnotesize\sffamily
  \yourname\quad
  web: \textcolor{blue}{\itshape\yourweb}\quad
  \textcolor{blue}{\youremail}}

\newcommand{\soptitle}{Statement of Purpose}
\newcommand{\yourname}{Hrishikesh Belagali}
\newcommand{\youremail}{belagal1@msu.edu}
\newcommand{\yourweb}{https://lonelyneutrin0.github.io/portfolio/}

\newcommand{\statement}[1]{\par\medskip
  \underline{\textcolor{blue}{\textbf{#1:}}}\space
}

\usepackage[
  colorlinks,
  breaklinks,
  pdftitle={\yourname - \soptitle},
  pdfauthor={\yourname},
  unicode
]{hyperref}

\begin{document}

\begin{center}\LARGE\soptitle\\
\large of \yourname\
\end{center}

\hrule
\vspace{1pt}
\hrule height 1pt

\bigskip
My fascination with quantum algorithms began with my experience in molecular dynamics 
research, in simulating biomolecular systems using classical computers. I have 
witnessed the applications and limitations of classical computing in this domain, 
which sparked my interest in quantum computing as a potential solution to these challenges. Through the 
$\langle \rangle$ program, I hope to deepen my understanding of quantum algorithms
and their applications in simulation. This aligns with my long-term goal of 
pursuing graduate school in quantum algorithms and contributing to research in the field. \\

As an undergraduate research assistant in the MSU-DOE Plant Research Laboratory, 
I have worked on several projects in computational biophysics using molecular dynamics with VMD and NAMD.
These covered a broad range of membrane-protein systems, including thylakoid membrane-protein interactions, aquaporin capping 
mechanisms and the effect of small molecules on bacterial membrane dynamics. My experience involved setting up simluation systems, 
running simulations on high performance computing clusters using SLURM, and analyzing the results using Python and TCL scripts 
with packages such as NumPy, Scipy, and MDAnalysis. I've learnt that atomistic simulations can provide invaluable insights into 
biomolecular mechanisms that are difficult to obtain experimentally. However, these simulations are limited by 
classical computing power, which hampers the timescales and system sizes that can be studied. The limitations of classical 
simulation inspired my interest in quantum algorithms as a tool to expand the capabilities of molecular simulations. \\

As a Honors Research Scholar part of MSU-Q, I have been working on generalized subgroup simulators for quantum circuits 
under the guidance of Professor Ryan LaRose. I investigate computational group theory and compilation optimization methods 
to develop efficient quantum circuit simulators, using Python and Cirq. Such simulators are beneficial for tapping into the 
advantages of quantum computing in near-team applications. This experience has strengthened my programming skills 
and enhanced my understanding of quantum computing concepts. \\

In addition to my research experience, I have built projects in computational physics and quantum computing.
I am interested in Monte Carlo methods for optimization. I implemented a simulated annealing 
algorithm for optimization Quadratic Unconstrained Binary Optimization (QUBO) problems. I was curious 
about the potential of quantum computing to solve this problems more efficiently, 
so I built a quantum optimization algorithm which formulates the QUBO problem as an Ising model and 
simulates its adiabatic evolution using Python and Qiskit. The main challenge I faced in this project was making sure the evolution was slow enough to not leave the ground state of the 
system, which is crucial for maintaining accuracy. Additionally, I had to explore alternate computational avenues 
for modeling larger Ising models due to their high resource demands.
I have also had the opportunity to work on aqua-blue, an open-source package that implements echo state networks. 
This class of machine learning models is known for its high accuracy rates in predicting the evolution of chaotic systems 
with minimal computational demand. I implemented time lag handling and hyperparameter tuning for the model. 
I was able to use these models to predict the evolution of several dynamical systems, including the 
Lorenz system, Mackey-Glass system and the time evolution of a Gaussian wave packet in a quantum well.
By doing so, I became familiar with systems exhibiting spatiotemporal chaos and gained experience with scientific computing packages like NumPy and PyTorch,
as well as familiarity with version control software and modular software design. 
\end{document}
